\section{Reproducible numerical checks}
The analytical results were checked with the accompanying NumPy scripts.  The numerical work
is used only to verify formulas, singular values and asymptotic trends; no figure represents
hardware data or a quantum-runtime benchmark.

\begin{table}[htbp]
\centering
\caption{Two-level model with $g=1$ and the condition-number-minimizing metric.  Values are rounded only for display.}\label{tab:values}
\begin{tabular}{@{}r r r r r@{}}
\toprule
$\gamma$ & $\kappa_2(\eta_*)$ & $\|H^{-1}\|_2$ &
$\|h_*^{-1}\|_2$ & $r_{\mathrm{sing}}(H)$\\
\midrule
0.00 & 1 & 1 & 1.000 & 1.00\\
0.50 & 3 & 2 & 1.155 & 0.50\\
0.80 & 9 & 5 & 1.667 & 0.20\\
0.90 & 19 & 10 & 2.294 & 0.10\\
0.99 & 199 & 100 & 7.089 & 0.01\\
\bottomrule
\end{tabular}
\end{table}

\begin{figure}[htbp]
\centering
\includegraphics[width=.82\linewidth]{figures/conditioning.pdf}
\caption{Two-level benchmark: the Euclidean condition number of the original Hamiltonian
coincides with that of the minimizing metric for $g=1$, while the Hermitian representative
has spectral condition number one for $\gamma<1$.}\label{fig:conditioning}
\end{figure}

\begin{figure}[htbp]
\centering
\includegraphics[width=.82\linewidth]{figures/inversion_and_radius.pdf}
\caption{Inverse norms and exact Euclidean distance to singularity in the two-level model.
Here the exceptional-point limit and inversion singularity coincide, so the metric-based
certificate is sharp.}\label{fig:radius}
\end{figure}

\begin{figure}[htbp]
\centering
\includegraphics[width=.84\linewidth]{figures/three_level_metric_sandwich.pdf}
\caption{Optimal-metric asymptotics for the coupled three-level family.  The rigorous lower
bound $8/\varepsilon^2$ and the explicit asymptotically optimal metric have the same leading
constant, while the simpler constructive metric supplies a uniform finite-$\varepsilon$ upper
bound.}\label{fig:metric_sandwich}
\end{figure}

\begin{figure}[htbp]
\centering
\includegraphics[width=.84\linewidth]{figures/three_level_certificate_v8.pdf}
\caption{Separation between metric conditioning and Euclidean invertibility.  The exact
singularity radius of $A_\varepsilon$ remains bounded away from zero.  The rigorous upper bound
$\varepsilon/(2\sqrt2)$ on the best metric certificate and the certificate generated by the
asymptotically optimal feasible metric both vanish linearly and become asymptotically equal.}
\label{fig:threelevel}
\end{figure}

\begin{figure}[htbp]
\centering
\includegraphics[width=.86\linewidth]{figures/growing_chain_radius.pdf}
\caption{Direct singular-value calculations for the growing coupled chain at representative
dimensions.  The Euclidean singularity radius remains uniformly positive, while the optimized
metric certificate is bounded above by $\varepsilon/(2\sqrt2)$ independently of dimension.}
\label{fig:growingchainradius}
\end{figure}
